% ======================================================================
% INSTRUCTIONS WHEN EDITING THIS PROJECT (inherited from instrument/)
% ======================================================================
%
% 1. Do not guess labels or numbers. Always use \ref{...} exactly as
%    requested. Never invent labels.
% 2. Exact hierarchy: \chapter / \section / \subsection / \subsubsection.
% 3. ASCII only in body text. No Unicode punctuation.
% 4. SOFT MARGINS OF 100 characters for all generated prose.
% 5. NO Lean in the prose: no identifiers, no syntax, no tactic talk.
%    The carriers become mathematical statements. The apparatus may be
%    named once, in the bibliography, as apparatus.
% 6. Null-base arrows fix exposition direction: QED -> GR,
%    lambda-calculus -> BNF, Topology -> Algebra -> Geometry.
%    Sections land on the right-hand (structural/geometric) side.
%
% ======================================================================

\usepackage[margin=1.2in]{geometry}
\usetikzlibrary{positioning}

% Theorem environments
\declaretheorem[numberwithin=chapter]{theorem}
\declaretheorem[sibling=theorem]{lemma}
\declaretheorem[sibling=theorem]{claim}
\declaretheorem[sibling=theorem]{corollary}
\declaretheorem[style=definition,sibling=theorem]{definition}
\declaretheorem[style=definition,sibling=theorem]{experiment}
\declaretheorem[style=remark,sibling=theorem]{remark}
\declaretheorem[style=remark,sibling=theorem]{fence}

% The five reading-registers of a claim (from the outline's typography;
% rendered as marginal disciplines, never as tags in the prose).
\newcommand{\bracket}[2]{\left[\,#1,\ #2\,\right]}

\setstretch{1.08}
